对实状态 $a,b$，由乘积法则和对称性可知通量 $a^{\mathsf t}C^\mu b$ 的散度为 $(\mathcal P a)^{\mathsf t}b+a^{\mathsf t}\mathcal P b$。这正是 \cite{r-25} 的双线性守恒恒等式。在 $(w_\ell,\delta w_\ell)$ 处，其值为 $q^\mu$ 沿 $\delta w_\ell$ 方向的导数。在管坐标 $\mathsf X-\mathrm i q$ 中（参见 \cite{r-6}, Section~1），直接求导可见 $D_{\ell a}$ 作用于 $\mathsf X-\mathrm i q$ 的每个分量均得零。故 $D_{\ell a}\overline G=0$，取共轭得 $\partial_{w_{\ell a}}G=\mathrm i V_{\ell a}G$。每个 $D_{\ell a}$ 的散度均为零，因为它的状态系数为常数，且空间系数与 $\mathsf X$ 无关。因此
\[
\operatorname{div}\mathcal Z_\ell=2\mathrm i\sum_aV_{\ell a}G_a
=2\mathrm i w_\ell^{\mathsf t}\mathcal P G=0.
\]
其中 $\mathcal P G=0$ 由沿 $\mathsf X$ 求导并取共轭得到，因为所有矩阵均为实矩阵。

对通常的平直联络，$D_{\ell a}\overline G=0$ 使下列计算中只有 $D_{\ell a}$ 作用于 $\overline{D_{mb}}$ 的系数：
\[
\nabla_{\mathcal Z_\ell}\overline{\mathcal Z_m}
=-\mathrm i\delta_{\ell m}\sum_\mu(G^{\mathsf t}C^\mu\overline G)\partial_{\mathsf X^\mu}.
\]
事实上，系数 $(C^\mu w_m)_b$ 沿 $w_{\ell a}$ 的导数为 $\delta_{\ell m}(C^\mu)_{ba}$。与 $-\mathrm i\delta_{\ell m}$ 相乘的向量由 $C^\mu$ 的对称性可知为实向量。由乘积法则得 $\operatorname{div}(\mathcal Z_\ell\otimes\overline{\mathcal Z_\ell})=(\operatorname{div}\mathcal Z_\ell)\overline{\mathcal Z_\ell}+\nabla_{\mathcal Z_\ell}\overline{\mathcal Z_\ell}$，其实部为零。用实向量场写出 $\mathcal Z_\ell=A_\ell+\mathrm iB_\ell$，可见取实部后被加项为 $A_\ell\otimes A_\ell+B_\ell\otimes B_\ell$。同一平直联络恒等式的虚部给出 $[A_\ell,B_\ell]=\sum_\mu(G^{\mathsf t}C^\mu\overline G)\partial_{\mathsf X^\mu}$。
